\subsection{The same-row inverse fourth moment and its endpoint payment}\label{proof:Q5-C100}
\paragraph{Q100.2 Kernel notation and inherited estimates}

For finite trees with at most five branch vertices, the orientation
problem requires the following bound on every finite probability space
and every real skew \(|K|\le 1\):

\fitdisplay{
Z_D=E_{\mu^{V(U)}}\prod_{v\to w\in D}(1+K(x_v,x_w))\le1.
}{}



Delete zero-mass atoms only. On the unchanged original L2(mu), use

\fitdisplay{
Jv=(Ev)1,\quad A=J+K,\quad T=J-K=A^*,\quad
 g=K1,\quad p=1+g,\quad q=1-g,
}{}
\fitdisplay{
f_n=A^{n}\one,\quad \bar f_n=T^{n}\one,\quad h_n=T f_{n-1},\quad
 W_n=A \bar f_{n-1},\quad \delta_j=\|K^{j}\one\|^2,\quad d=\delta_1.
}{}

The finite calculus gives \(\theta=13/32,\allowbreak  \kappa=19/32,\allowbreak  d\le1/3\), \(d^2\le\delta_2\le\theta d\), \(\delta_{j+1}\le\theta\delta_j\), and \(\|K\|^2\le\theta\).
Skewness gives \(K(x,\allowbreak x)=0\), so p(x), \(q(x)\ge \mu(x)>0\) on retained atoms.
Every actual rooted tree message is positive by the all-equal labeling.

The proof uses the finite calculus and the four-branch pairing and zero-step interfaces already established above.

The genuine equal-pair comparison is \(f_l h_l\le 1,\allowbreak \) or its full reverse
\(\bar f_l W_l\le 1.\) Each such factor stays at its original vertex and is removed
only from a nonnegative upper comparison. The path-mean bound used here applies to $f_n$ and $\bar f_n$. In contrast, $\E h_2=\E W_2=1+d$.
The f2 case uses the non-strict bound.

\paragraph{Q100.3 Independent finite-G reconstruction}

Put \(G_m(z)=\sum _{j=0}^{m-1} z^j\) and
\(E_r(x)=\sum _{j=0}^{r-1}(-x)^j.\) For \(0\le x\le \theta<1,\allowbreak \) the finite
alternating sum satisfies \(1-x\le E_r\le 1.\) At \(z=i \sqrt{x}\), write
\(G_m(z)=A_m(x)+zD_m(x).\) For the actual

\fitdisplay{
R=K G_m(-K)g,\qquad h=E(gR),\qquad V=ER^2,
 \qquad m=n-1\ge3,
}{}

the block expression is \(R=(zA_m+xD_m)g.\) Hence the genuine finite
spectral measure \(\nu_g\) of -\(K^2\) at g gives

\fitdisplay{
h=\int xD_m\,d\nu_g,\qquad
 V=\int x(A_m^2+xD_m^2)\,d\nu_g.
}{}

Its mass is d and its first moment is \(\delta_2.\) This proves the signs and
identifies the same R as in the original correlated cubic calculation.

\paragraph{Q100.3.1 Even and odd finite lengths}

For \(m=2r,\allowbreak  r\ge 2,\allowbreak  A_m=D_m=E_r,\allowbreak \) and the FINITE identity
(\(1+x)E_r=1-(-x)^r\) yields

\fitdisplay{
x[(1+x)E_r^2-E_r]=(-1)^{r+1}x^{r+1}E_r.
}{}

Positive cases have r odd, so \(r\ge 3.\) They are bounded by \(x^4,\allowbreak \)
because \(E_r\le 1\) and \(x\le 1.\) Furthermore
\(x^4\le x^3(1-x+x^2),\allowbreak \) since their difference is \(x^3(1-x)^2.\)
Negative cases automatically satisfy the same upper bound.

For \(m=2r+1,\allowbreak  r\ge 1,\allowbreak  D_m=E_r\) and \(A_m=1-xE_r.\) Substitution, or replacing
(\(1+x)E_r\) by its finite expression, gives

\fitdisplay{
A_m^2+xD_m^2-D_m=(-x)^r A_m.
}{}

Here \(A_m>0.\) Only even \(r\ge 2\) gives a positive term. In that case

\fitdisplay{
A_m=\frac{1+x^{r+1}}{1+x}\le
 \frac{1+x^3}{1+x}=1-x+x^2,\qquad x^{r+1}\le x^3.
}{}

Multiplication by x proves the same uniform upper polynomial. The two
short endpoints are explicit:

\fitdisplay{
m=3:\ -x^2(1-x),\qquad
 m=4:\ -x^3(1-x).
}{}

They are negative, not missing positive cases. At \(m=5\) the expression
is exactly \(x^3(1-x+x^2),\allowbreak \) so that polynomial is a real finite expression.
These calculations cover every integer \(m\ge 3,\allowbreak \) with no infinite-tail
substitution.

\paragraph{Q100.3.2 Uniform spectral payment and exact constants}

The derivative of \(x^2(1-x+x^2)\) is \(x(2-3x+4x^2),\allowbreak \) nonnegative;
the quadratic has discriminant -23 and positive leading coefficient.
Therefore

\fitdisplay{
\begin{aligned}
 V-h
 &\le\int x^3(1-x+x^2)\,d\nu_g\\
 &\le\theta^2(1-\theta+\theta^2)\int x\,d\nu_g
 =r_*\delta_2.
 \end{aligned}
}{}

Since \(1-\theta+\theta^2=777/1024\) and \(\theta^2=169/1024,\allowbreak  r_*=131313/1048576.\) The comparisons in (2.6)-(2.7) check exactly:

\begin{itemize}
\item \(r_*>1/8\) because \(131313>131072.\)
\item \(r_*<13/100\) because \(13131300<13631488.\)
\item \(\theta r_*=1707069/33554432<51/1000\) because
  \(1707069000<1711276032.\)
\item \(\theta^2-51/1000=(21125-6528)/128000=14597/128000.\)
  This is >\(9/80\) because \(14597>14400.\)
\end{itemize}

Thus \(\theta(\theta-r_*)>\delta_*=14597/128000>0.\) No step assumes \(x\ge d.\)
In particular this is a correlated V-h bound, not independent worst-case
budgets for V and h.

\paragraph{Q100.4 Exact same-row calculation}

This section reconstructs the calculation through (2.10). Write \(\alpha=\kappa d\), \(C=1-\alpha,\allowbreak  t=1+d/2\) and \(\beta=\beta_{n+1}.\)

The finite convolution is

\fitdisplay{
S'_n=\sum_{j=0}^n(-K)^j1
 =\bar f_n+\sum_{j=3}^n\beta_j\bar f_{n-j}=q+R.
}{}

For \(n\ge 4\) retain its actual b0 and b1 terms. Since \(\beta_n,\allowbreak  \beta_{n-1}\)
are at least alpha, \(\bar f_n\le Cq+R-\alpha.\) Also T \(v\le 2Ev\) and
\(T^2v\le 2(Ev)q\) for \(v\ge 0.\) The same finite convolution gives
\(E\bar f_{n-2}\le 1-\beta_{n-1}\le C,\allowbreak \) hence \(\bar f_n\le 2Cq.\)
At \(n=4\) the latter uses the actual b2 mean 1-d; no short term is lost.

For \(s\ge -1,\allowbreak \)

\fitdisplay{
\min(1+s,2)^4\le1+4s+7s^2+4s^3.
}{}

On [-1, 1] the difference from (\(1+s)^4\) is \(s^2(1-s^2).\)
On [1, infinity), the cubic starts at 16 and has positive derivative.
The upper Cq+R-alpha is positive, so \(s=(R-\alpha)/(Cq)>-1.\)
Applying the inequality to the same actual row gives

\fitdisplay{
B_n\le C^4Eq^3+4C^3E[q^2(R-\alpha)]
       +7C^2E[q(R-\alpha)^2]+4CE(R-\alpha)^3.
}{Q100.A24.1}

The exact finite quantities satisfy \(ER=-\beta,\allowbreak  \alpha\le \beta\le d,\allowbreak  \kappa\delta_2\le h\le\delta_2\), \(V\le \delta_2,\allowbreak  |R|\le \sqrt{d},\allowbreak \) and
E[\(G_m(-K)g\)]=\(\beta_n\ge \alpha.\) The h and V bounds follow from the
actual finite coordinates \(D_m\ge 1-x\) and \(|G_m|\le 1\) for \(m\ge 3.\)

Put \(W=2R-Cg.\) The original row estimate gives

\fitdisplay{
|W|^2\le C^2+4d-4C\alpha
 =1+(4-6\kappa)d+5\kappa^2d^2\le t^2.
}{}

The last difference is
\(d[9/16-(1549/1024)d]\ge179d/3072\). Also

\fitdisplay{
EW=-2\beta,\qquad D_W=EW^2=C^2d+4V-4Ch\ge0.
}{}

The pointwise inequality (\(W+2\beta)^2(W-t)\le 0\) gives

\fitdisplay{
EW^3\le(t-4\beta)D_W-4t\beta^2+8\beta^3.
}{}

The full correlated identity is

\fitdisplay{
-C^4Eg^3+4C^3E(g^2R)-6C^2E(gR^2)+4CER^3
 =\frac{C}{2}[EW^3+C^2E(g^2W)].
}{}

Use \(E(g^2W)\le td.\) Of the seven copies of -\(C^2E(gR^2)\) in the
expansion (Q100.A24.1), six are in this identity and one remains. This is
where a separate absolute cubic budget would lose the intended gain.

At fixed TRUE \(D_W\ge 0,\allowbreak \) the beta-linear coefficient from the rest of
(Q100.A24.1) is

\fitdisplay{
-4C^3+14C^2\alpha-12C\alpha^2
 =-2C(2C-3\alpha)(C-2\alpha)<0.
}{}

Indeed \(\alpha\le 19/96<1/5.\) The derivative of the cubic bound, after
its C/2 factor, is
\(C[-2D_W-4t\beta+12\beta^2]\le0\) because \(3\beta\le 1\le t.\)
Thus one first lowers beta to alpha, while holding \(D_W\) fixed.
Only afterwards \(t-4\alpha=1-15d/8\ge 3/8\) allows its substitution.

Before grouping V-h, the coefficients are precisely

\fitdisplay{
C[7+2t-27\alpha]V+C^2[-8-2t+30\alpha]h.
}{}

Writing \(V=(V-h)+h\) makes the h coefficient

\fitdisplay{
\begin{aligned}
 C\{7+2t-27\alpha+(1-\alpha)(-8-2t+30\alpha)\}
 &=C[-1+(11+2t)\alpha-30\alpha^2]\\
 &=C[-1+(13+d)\alpha-30\alpha^2].
 \end{aligned}
}{Q100.A24.2}

This independently proves (2.10), including the disputed
coefficient. Together with its unchanged constants it yields

\fitdisplay{
\begin{aligned}
 B_n\le{}&C^4(1+3d)+C^3[td-6\alpha d-8\alpha]
             +C\alpha^2[19-d-33\alpha]\\
 &+CB(d)(V-h)
 +C[-1+(13+d)\alpha-30\alpha^2]h-C^2E(gR^2),
 \end{aligned}
}{Q100.A24.3}

where \(B(d)=7+2t-27\alpha=9-481d/32\) and \(B\ge 383/96>0.\)


\paragraph{Q100.4.2 Complete signed substitution}

Split the actual h coefficient
term by term, using the lower h bound on its negative terms and the
upper h bound on its positive term:

\fitdisplay{
[-1+(13+d)\alpha-30\alpha^2]h
 \le[-\kappa+13\kappa d+\kappa(1-30\kappa^2)d^2]\delta_2
 =H(d)\delta_2.
}{}

After the proved V-h improvement, the combined coefficient is
\(S_new=r_*B+H.\) Its constant is >\(17/32,\allowbreak \) using \(r_*>1/8.\)
Its linear coefficient is

\fitdisplay{
13\kappa-\frac{481}{32}r_*
 >\frac{24700-6253}{3200}
 =\frac{18447}{3200}>5.
}{}

Its quadratic coefficient is -\(93157/16384>-6.\) Thus

\fitdisplay{
S_{\rm new}>17/32+5d-6d^2>0,\qquad 0\le d\le1/3.
}{}

This sign is proved BEFORE multiplying the upper bound \(\delta_2\le\theta d\).
\paragraph{Q100.5 Old polynomial, exact difference and continuous row}

For the coefficient P, let F0 be
the first three terms of (Q100.A24.3). The combined coefficient is

\fitdisplay{
S_{\rm old}=\theta B+H
 =\frac{49}{16}+\frac{1651}{1024}d
       -\frac{93157}{16384}d^2.
}{}

For example, its linear coefficient is
-\(6253/1024+7904/1024=1651/1024.\) 

The old polynomial is \(P=F_0+C\theta d S_{\rm old}\). Its coefficient ledger in kappa is


\begingroup\small
\setlength{\tabcolsep}{3pt}
\setlength{\LTleft}{0pt}\setlength{\LTright}{\fill}
\begin{longtable}{@{}>{\raggedright\arraybackslash}p{0.1200\linewidth}>{\raggedright\arraybackslash}p{0.4000\linewidth}>{\raggedright\arraybackslash}p{0.4000\linewidth}@{}}
\caption{Coefficients of \(d^j\) in the comparison polynomial \(P(d)\).}\label{tab:parameters:q5-c100:1}\\
\toprule
\(j\) & Coefficient before setting \(\kappa=19/32\) & Exact value \\
\midrule\endfirsthead
\toprule
\(j\) & Coefficient before setting \(\kappa=19/32\) & Exact value \\
\midrule\endhead
\bottomrule\endfoot
1 & \(13-31\kappa+10\kappa^2\); & \(-963/512\). \\
\addlinespace[2pt]
2 & \(3/2-46\kappa+110\kappa^2-37\kappa^3\); & \(171113/32768\). \\
\addlinespace[2pt]
3 & \(-(3/2)\kappa+53\kappa^2-152\kappa^3+57\kappa^4\); & \(-7275575/1048576\). \\
\addlinespace[2pt]
4 & \((1/2)\kappa^2-29\kappa^3+72\kappa^4-30\kappa^5\); & \(14103187/16777216\). \\
\addlinespace[2pt]
5 & \(9\kappa^4-(1/2)\kappa^3\); & \(1063145/1048576\). \\
\end{longtable}
\endgroup


Expanding the finite products gives these coefficients. For instance the degree-two numerator
is \(49152-894976+\;1270720-253783=171113\); the degree-three numerator is
-\(933888+\;19592192-33362176+\;7428297=-7275575.\) At degree four, the
numerator over 33554432 is
\(5914624-203684864+\;300259584-74282970=28206374,\allowbreak \) giving the printed
fraction after division by two. At degree five it is
\(1172889-109744=1063145\) over 1048576.

The correct new polynomial is consequently EXACTLY

\fitdisplay{
F0+C\theta d(r_*B+H)
 =P-CB\theta(\theta-r_*)d.
}{Q100.A24.4}

The correlation satisfies

\fitdisplay{
|E(gR^2)|\le\sqrt{d\,ER^4}
 \le d\sqrt V\le\sqrt\theta\,d^{3/2}\le\frac23d^{3/2}.
}{}

Here \(|R|^2\le d\) and \(V\le \theta\) d are the actual same-R bounds.
Combining them with (Q100.A24.4) proves (3.4) exactly.

The rational majorization of P also checks:
\(A2<21/4\) follows from \(171113<172032\);
\(A3<-69/10\) from \(72755750>72351744\);
\(A4<1\) and \(A5<33/32\) from
\(14103187<16777216\) and \(1063145<1081344.\)
Since

\fitdisplay{
d^3+\frac{33}{32}d^4
 \le\left(\frac13+\frac{33}{288}\right)d^2
 =\frac{43}{96}d^2,
}{}

and \(69/10-43/96=3097/480>32/5,\allowbreak \) one gets (3.5) for \(d>0.\)

Using \(\theta(\theta-r_*)>\delta_*\) proves

\fitdisplay{
B_n\le1-s(d)d,
}{}
\fitdisplay{
s(d)=\frac{963}{512}-\frac{21}{4}d+\frac{32}{5}d^2
       -\frac23(1-\kappa d)^2\sqrt d
       +\delta_*(1-\kappa d)(9-481d/32).
}{}

All factors in its last product are positive. The derivative of its first
three terms is at most -\(21/4+64/15=-59/60.\) The product
(1-kappa d)(9-481d/32) decreases, since both positive factors decrease.
For \(d>0,\allowbreak \)

\fitdisplay{
\frac{d}{dd}\{(1-\kappa d)^2\sqrt d\}
 =\frac{(1-\kappa d)(1-5\kappa d)}{2\sqrt d}>0;
}{}

at \(d=1/3,\allowbreak \) 1-5kappa \(d=1/96.\) The negative of this term therefore
decreases too. These exact signs prove decreasing s on the whole
interval, with its continuous extension at zero. 

\paragraph{Q100.6 Completing the inverse-row bound.}
The final scalar payment for
\(B_n=\E[\bar f_n^4/q]\), \(n\ge4\), is as follows.
The derivative above is positive on \(0<d\le1/3\), and
\(1/\sqrt3<11/19\), since \(361<363\). Hence
\fitdisplay{
\frac23(1-\kappa d)^2\sqrt d
 \le\frac23\left(\frac{77}{96}\right)^2\frac{11}{19}
 =\frac{130438}{525312}<\frac14.
}{Q100.3.6}
The last comparison is \(521752<525312\).
Dropping the nonnegative improvement in \(s(d)\), the preceding
coefficient bound therefore gives, for \(d>0\),
\fitdisplay{
\frac{B_n-1}{d}
 \le-\frac{835}{512}+\frac{21}{4}d-\frac{32}{5}d^2.
}{Q100.3.7}
This quadratic increases on \([0,\allowbreak 1/3]\): its derivative is at least
\(21/4-64/15=59/60>0\). Its value at the right endpoint is
\fitdisplay{
-\frac{835}{512}+\frac74-\frac{32}{45}
 =-\frac{13639}{23040}<-\frac12.
}{Q100.3.8}
At \(d=0\), \(\bar f_n=q=1\). Thus every finite \(n\ge4\) satisfies
\fitdisplay{
\boxed{\E[\bar f_n^4/q]\le1-d/2.}
}{Q100.3.9}
The separate \(n=3\) estimate in Section~\ref{proof:Q5-C056} is
\(1-23d/36\); since \(23/36>13639/23040>1/2\), it implies the
same weaker bound. These are the actual inward paths for every
length throughout the score interval.
